\subsection{معرفی و انگیزه}
این مکانیزم بر مبنای مبانی نظری پیشرفته در ادبیات توجه و فشرده‌سازی \cite{ref_roasthashedlinear_2024} توسعه یافته است.

اشتراک‌گذاری پارامترها با استفاده از توابع درهم‌سازی تصادفی (\lr{Random Hashing}) رویکردی کارآمد برای کاهش ردپای حافظه ماتریس‌های وزن بدون نیاز به نگهداری ماتریس‌های کامل در فضای رم به شمار می‌رود. با این وجود، روش‌های سنتی مانند ترفند درهم‌سازی اسکالر (\lr{Scalar Hashing}) به دلیل عدم تطابق با الگوهای دسترسی حافظه در پردازنده‌های گرافیکی و نقض محلی بودن داده‌ها (\lr{Spatial Locality})، منجر به ناکارآمدی حافظه نهان و افت فاحش سرعت اجرای مدل می‌گردند.

ماژول \lr{roast\_hashed\_linear} با نام فارسی «لایه خطی درهم‌سازی‌شده با کاشی‌بندی ROAST» مبتنی بر ایده مقاله مرجع \lr{ROAST: A Hardware-Aware Weight Sharing Scheme for Efficient Neural Networks} توسعه یافته است. در این الگوریتم، وزن‌های شبکه به جای عناصر اسکالر منفرد، در سطح بلوک‌ها یا کاشی‌های دوبعدی (\lr{Tiling}) با جدول کوچکی از مقادیر به اشتراک گذاشته می‌شوند که سازگاری کاملی با سلسله‌مراتب حافظه سخت‌افزاری ایجاد می‌نماید.
